% Beamer template for academic finance and real-estate presentations
% Copy this file and rename for your presentation
% Compile: pdflatex -> pdflatex (no bibtex needed unless using citations)

\documentclass[aspectratio=169,11pt]{beamer}

% Theme
\usetheme{metropolis}
\usepackage{appendixnumberbeamer}

% Colors
\definecolor{NavyBlue}{RGB}{0,40,85}
\definecolor{DarkGray}{RGB}{64,64,64}
\definecolor{MedGray}{RGB}{128,128,128}
\definecolor{LightGray}{RGB}{220,220,220}
\definecolor{AccentRed}{RGB}{180,30,30}

\setbeamercolor{frametitle}{bg=NavyBlue}
\setbeamercolor{progress bar}{fg=NavyBlue}
\setbeamercolor{alerted text}{fg=AccentRed}

% Packages
\usepackage{booktabs}
\usepackage{graphicx}
\usepackage{amsmath,amssymb}
\usepackage{tikz}
\usepackage{hyperref}
\usepackage[natbib=true,style=authoryear-comp,maxcitenames=2,giveninits=true,uniquename=false,uniquelist=false,doi=false,isbn=false,url=false]{biblatex}
% \addbibresource{references.bib}  % Uncomment and set path to your .bib file

% Custom commands
\newcommand{\sym}[1]{\ensuremath{^{#1}}}
\newcommand{\backupbegin}{\newcounter{framenumberappendix}\setcounter{framenumberappendix}{\value{framenumber}}}
\newcommand{\backupend}{\addtocounter{framenumberappendix}{-\value{framenumber}}\addtocounter{framenumber}{\value{framenumberappendix}}}

% Settings
\setbeamertemplate{navigation symbols}{}  % Remove navigation buttons
% \setbeameroption{show notes on second screen=right}  % Uncomment for dual-screen
% \pgfpagesuselayout{2 on 1}[a4paper]  % Uncomment for handout layout

% Metadata
\title{Paper Title}
\subtitle{Subtitle or Conference Name}
\author{Author Name \\ \small{Affiliation}}
\date{\today}
\institute{}

\begin{document}

% ============================================================
% TITLE
% ============================================================
\begin{frame}[plain]
\titlepage
\end{frame}

\begin{frame}{Contribution}
\begin{center}
\large
[One-sentence contribution statement]
\end{center}
\note{State the contribution clearly in the first 30 seconds.}
\end{frame}

% ============================================================
% MOTIVATION
% ============================================================
\begin{frame}{Motivation}
\begin{itemize}
    \item [Recent event, statistic, or trend that motivates the question]
    \item [Why this matters economically]
    \item [The gap or tension this creates]
\end{itemize}
\vfill
\textcolor{MedGray}{\scriptsize Source: [citation or data source]}
\note{Connect to something the audience already knows or cares about.}
\end{frame}

% ============================================================
% RESEARCH QUESTION
% ============================================================
\begin{frame}{Research question}
\begin{center}
\large
\textbf{[Precise, testable question]}
\end{center}
\vspace{1cm}
\begin{itemize}
    \item Prediction 1: [If mechanism A, then...]
    \item Prediction 2: [If mechanism B, then...]
\end{itemize}
\note{The audience should be able to hold this question in their head for the entire talk.}
\end{frame}

% ============================================================
% INSTITUTIONAL SETTING / MECHANISM
% ============================================================
\begin{frame}{Institutional setting}
\begin{itemize}
    \item [Key institutional feature 1]
    \item [Key institutional feature 2]
    \item [Why this creates useful variation for identification]
\end{itemize}
\note{Build the audience's intuition before showing the design. Do not rush this slide.}
\end{frame}

% ============================================================
% PREVIEW (required for conferences, optional for seminars)
% ============================================================
\begin{frame}{Preview of results}
\begin{enumerate}
    \item [Main finding with magnitude]
    \item [Key secondary finding]
    \item [Mechanism evidence summary]
\end{enumerate}
\note{Conferences need this early. Seminars can defer to build suspense.}
\end{frame}

% ============================================================
% IDENTIFICATION
% ============================================================
\begin{frame}{Identification challenge}
\begin{itemize}
    \item [What makes this question hard to answer?]
    \item [What would a naive estimate get wrong?]
    \item [The specific endogeneity concern]
\end{itemize}
\note{Acknowledge the challenge before presenting the solution. Builds credibility.}
\end{frame}

\begin{frame}{Research design}
\begin{equation*}
    \textit{Outcome}_{i,t} = \beta \, \textit{Treatment}_{i,t} + \gamma' X_{i,t} + \alpha_i + \delta_t + \varepsilon_{i,t}
\end{equation*}
\vspace{0.5cm}
\begin{itemize}
    \item $\textit{Treatment}_{i,t}$: [definition]
    \item Fixed effects: [what they absorb]
    \item Clustering: [level]
    \item Identifying assumption: [in plain language]
\end{itemize}
\note{This is the slide where skeptics decide whether to trust you.}
\end{frame}

% ============================================================
% DATA
% ============================================================
\begin{frame}{Data and sample}
\begin{itemize}
    \item Sources: [list]
    \item Sample period: [dates]
    \item Observations: [N firms/properties/borrowers $\times$ T periods]
    \item Key variables: [treatment, outcome, controls]
\end{itemize}
\note{Keep brief. One slide unless the data construction is a contribution.}
\end{frame}

% ============================================================
% MAIN RESULTS
% ============================================================
\begin{frame}{[Economic finding, not ``Main results'']}
\begin{center}
    % Prefer coefficient plots or event-study graphs over tables
    % \includegraphics[width=0.85\textwidth]{output/figures/fig_main_result.pdf}
    % TODO: Replace with actual figure path
    \textit{[Coefficient plot or event-study figure here]}
\end{center}
\vspace{-0.5cm}
{\scriptsize Controls: [list]. FE: [list]. SE clustered at [level].}
\note{Show the result, then explain the magnitude.}
\end{frame}

% ============================================================
% MAGNITUDE
% ============================================================
\begin{frame}{Economic magnitude}
\begin{center}
\large
A one-SD increase in \textit{Treatment} \\[0.5cm]
$\Rightarrow$ \textbf{[Y]\% change} in \textit{Outcome} \\[0.5cm]
$\approx$ \textbf{\$[amount] per [unit]} \\[0.5cm]
{\normalsize (equivalent to [familiar benchmark])}
\end{center}
\note{Three tiers: statistical units, real-world units, anchored comparison.}
\end{frame}

% ============================================================
% MECHANISM
% ============================================================
\begin{frame}{Mechanism: [channel name]}
\begin{itemize}
    \item If [channel A]: expect [pattern 1]
    \item If [channel B]: expect [pattern 2]
    \item Evidence: [which pattern we observe and what it rules out]
\end{itemize}
% \includegraphics[width=0.7\textwidth]{output/figures/fig_mechanism.pdf}
\note{Be honest about what the mechanism evidence can and cannot show.}
\end{frame}

% ============================================================
% ROBUSTNESS (keep to 1-2 slides)
% ============================================================
\begin{frame}{Robustness}
\begin{itemize}
    \item \textbf{[Threat 1]}: [test and result]
    \item \textbf{[Threat 2]}: [test and result]
    \item \textbf{[Threat 3]}: [test and result]
\end{itemize}
\note{Name the threat, show the check, move on. Do not dwell.}
\end{frame}

% ============================================================
% CONCLUSION
% ============================================================
\begin{frame}{Takeaway}
\begin{enumerate}
    \item \textbf{Finding}: [main result with magnitude]
    \item \textbf{Mechanism}: [channel with supporting evidence]
    \item \textbf{Implication}: [economic lesson or policy relevance]
\end{enumerate}
\vfill
\begin{center}
{\large author@university.edu}
\end{center}
\note{Do not end with ``thank you.'' End with the idea.}
\end{frame}

% ============================================================
% BACKUP SLIDES
% ============================================================
\appendix
\backupbegin

\begin{frame}[noframenumbering]{Backup: Alternative specifications}
% [Content]
\end{frame}

\begin{frame}[noframenumbering]{Backup: Pretrend tests}
% [Content]
\end{frame}

\begin{frame}[noframenumbering]{Backup: Sample construction}
% [Content]
\end{frame}

\begin{frame}[noframenumbering]{Backup: Comparison to [Paper]}
% [Content]
\end{frame}

\backupend

\end{document}
